\section{Introduction}

The previous chapter described the AI-Scientist architecture as an orchestrated verification pipeline. This chapter
maps that design onto the actual implementation. The goal is not to list every file mechanically, but to explain
how the main modules cooperate to realize evidence-grounded, self-critical, traceable verification in code.

The implementation is organized as a single Python package, \texttt{ai\_scientist}, with the agents in an
\texttt{agents} subpackage and a set of supporting modules for configuration, data models, corpus access,
ingestion, retrieval utilities, scope control, and the optional language-model layer. A Streamlit user interface
and a programmatic API sit on top of the same package. This separation is intentional: the package contains the
reasoning logic, while the interfaces are thin clients that drive it.

\section{Implementation Structure}

At the highest level, the codebase is divided into:

\begin{enumerate}
\item \textbf{agent modules} under \path{ai_scientist/agents/}, which implement retrieval, claim extraction,
verification, critique, synthesis, reporting, and live retrieval;
\item \textbf{the orchestrator} in \path{ai_scientist/agents/orchestrator.py}, which wires the agents into the
end-to-end workflow;
\item \textbf{core modules} --- \path{models.py}, \path{config.py}, \path{corpus.py}, \path{ingestion.py},
\path{utils.py}, \path{scope.py}, and \path{llm.py} --- which provide the data types, settings, evidence access,
and shared logic;
\item \textbf{baselines} in \path{baselines.py}, which implement the single-agent and RAG comparison systems on
top of the same retrieval and scoring code; and
\item \textbf{interfaces}: a Streamlit application (\path{app.py}) and a FastAPI service
(\path{ai_scientist/api.py}).
\end{enumerate}

This layered structure keeps the reasoning system reusable. The same \texttt{analyze\_question(...)} entry point
used by the user interface is also used by the API and by the baselines harness, so that what is demonstrated is
the real system rather than a simplified mock-up.

\section{Typed Data Models}

The pipeline communicates through explicit dataclasses defined in \path{models.py}. A \texttt{PaperDocument}
carries a paper's identifier, title, abstract, year, authors, and topics. A \texttt{StructuredClaim} records a
claim's text, type, source paper, source sentence, and focus terms. An \texttt{EvidenceSnippet} pairs a sentence
with its overlap score and source. A \texttt{ClaimAssessment} bundles a claim with its verdict, confidence,
evidence list, and rationale. Finally, a \texttt{FinalReport} aggregates the question, summary, verified claims,
critique notes, retrieved papers, rendered Markdown, and the iteration trace.

These typed structures are important because every stage writes to a shared, serializable interface rather than
passing unstructured text. This is what makes the system auditable: the report is not a paragraph that happens to
mention sources, but a structured object in which each conclusion is explicitly linked to its evidence.

\section{The Orchestrator}

The practical center of the implementation is \texttt{MultiAgentResearchSystem} in \path{orchestrator.py}. Its
\texttt{analyze\_question} method realizes the workflow described in Chapter~4: it applies the scope guard,
retrieves evidence, extracts claims, verifies them through a refinement routine, runs the Critic, optionally
re-verifies weak claims, synthesizes the summary, and builds the report. The method also implements the honest
early exits: if the question is out of scope, if retrieval returns no papers, or if no aligned claims can be
extracted, it returns a report that explains the situation instead of fabricating an answer.

Adaptive re-verification is implemented in the private \texttt{\_refine\_assessments} and \texttt{\_run\_second\_pass}
routines. After the first pass, the Critic's \texttt{identify\_reverification\_targets} produces a list of
follow-up decisions; for each, the Orchestrator retrieves additional papers for that claim's follow-up query,
merges them with the first-pass evidence, and re-verifies the claim with contradiction resolution preferred. Every
step appends a human-readable line to the iteration trace, which is what the user ultimately sees as the
orchestration trace.

\section{Retrieval and the Layered Evidence Pool}

The Retrieval agent in \path{agents/retrieval.py} implements the layered evidence model. Its \texttt{\_local\_pool}
builds a deduplicated pool that lists uploaded papers before curated-corpus papers, and \texttt{\_local\_rank\_score}
adds the source-dependent bonus so that uploads outrank corpus papers at comparable relevance. The
\texttt{\_rank\_subset} routine ranks candidates by this source-weighted score but gates admission on the raw
query-coverage clearing the configured minimum, so the bonus can reorder evidence but never inject an irrelevant
paper.

When live retrieval is enabled, \texttt{retrieve} layers API results on top of the local pool. Live papers pass
through a quality filter that rewards adequate abstract length, query relevance in title and abstract, recency, and
the presence of research-indicator terms, and they are merged with a strict source priority: uploads first, then
corpus, then live. A title-similarity check removes near-duplicates so the same paper is not counted twice across
sources. The agent is wrapped so that any network failure falls back cleanly to the local pool, keeping the system
usable offline.

\section{Retrieval Utilities}

The relevance measures live in \path{utils.py}. Tokenization normalizes case, strips stopwords, and applies a
small morphological normalization so that, for example, plural and singular forms of common research terms match.
On top of tokenization, \texttt{overlap\_score} computes Jaccard overlap for short-text comparison, and
\texttt{coverage\_score} computes the length-robust query coverage used for paper-level ranking. Keeping both
measures in one module makes the design decision from Chapter~3 --- coverage for ranking, Jaccard for
sentence-level evidence --- explicit and reusable across the verifier, the retriever, and the baselines.

\section{Claim Extraction}

The Claim Extraction agent in \path{agents/claim_extraction.py} scans each retrieved paper sentence by sentence. A
sentence becomes a candidate claim if it contains a positive research cue or shares content tokens with the
question. Each candidate is wrapped in a \texttt{StructuredClaim} whose type is inferred from surface cues ---
limitation, performance, finding, implication, or observation --- and whose focus terms are the intersection of the
question and sentence tokens. Extraction stops once a configured maximum number of claims is reached, keeping the
workflow focused on the statements most relevant to the query.

\section{Verification}

The Verification agent in \path{agents/verification.py} is the analytical heart of the implementation. For each
claim it iterates over the sentences of the candidate papers and computes an evidence score that combines Jaccard
overlap with focus-term and content-term agreement and small type- and source-aware bonuses. Each scored sentence
is then classified as supporting or contradicting the claim, using negation and contradiction-cue analysis so that
a polarity mismatch between claim and sentence is treated as a contradiction.

From the collected evidence the agent selects a verdict. An overgeneralized claim returns insufficient evidence
unless a contradiction pass is explicitly requested; strong, multi-source counter-evidence yields a contradicted
verdict; adequate aligned support yields a supported verdict; and the absence of sufficient overlap yields
insufficient evidence. Confidence is computed by \texttt{\_supported\_confidence} and \texttt{\_contradicted\_confidence}
from the strongest evidence overlap, the count of additional supporting snippets, and the number of independent
corroborating sources, with an explicit penalty when a claim is corroborated only by its own source. All confidence
values are clamped between a floor and a ceiling so the system never reports absolute certainty. The optional
language-model refinement is invoked only after this deterministic assessment is formed, and only over the
evidence already selected.

\section{Critique and Adaptive Re-Verification}

The Critic agent in \path{agents/critic.py} implements both reporting and adaptivity. Its \texttt{critique} method
produces notes for insufficient, low-confidence, single-source, and contradicted claims, including details such as
whether a single-source claim is merely self-corroborated. Its \texttt{identify\_reverification\_targets} method
selects which weak claims warrant another pass and builds a follow-up query from the claim text, its focus terms,
its source title, and a counter-evidence probe when the claim uses absolute or contradiction-prone language.
Overgeneralized claims are excluded, because additional retrieval cannot settle a claim that is too broad. This is
the implementation of the bounded, targeted adaptivity required in Chapter~3.

\section{Synthesis and Reporting}

The Synthesis agent in \path{agents/synthesis.py} aggregates the assessments into a calibrated summary, optionally
rewritten by the language model under instructions to stay grounded and avoid overclaiming, with a deterministic
fallback. The Report agent in \path{agents/report.py} renders the structured \texttt{FinalReport}, including the
retrieved papers, the orchestration trace, and the full per-claim verification details with evidence snippets and
overlap scores. Because the report is both a structured object and rendered Markdown, the same result can be shown
in the interface, returned by the API, or exported.

\section{The Optional Language-Model Layer}

The language-model abstraction in \path{llm.py} wraps a provider client behind a single interface. It is available
only when an API key is configured; otherwise the agents run on their deterministic logic. The module is written to
be provider-tolerant: a helper detects whether the configured model accepts a reasoning-effort parameter and
includes it only for reasoning-capable models, while models that would reject it are called without it. Structured
outputs are parsed into small typed schemas for the refined verdict and the rewritten summary, and any call failure
is logged and degrades to the deterministic result rather than breaking the pipeline.

This design serves three purposes. First, it makes the core pipeline reproducible and offline-capable. Second, it
ensures the model can improve fluency and borderline verdicts but can never bypass the grounding and honesty rules,
which run first. Third, it keeps the system model-agnostic, so a different model or provider can be substituted
through configuration.

\section{Live Scholarly Retrieval}

Live evidence is implemented in \path{agents/live_retrieval.py}. The agent queries PubMed Central through the NCBI
E-utilities (search, summary, and fetch) and arXiv through its Atom API. Two implementation details are important
for evidence quality. First, PubMed Central searches are sorted by relevance rather than date, because the default
date ordering surfaces recent but off-topic papers. Second, abstracts in the PubMed Central XML are nested inside
structured elements, so the agent gathers all descendant text rather than only the top-level text node, which is
what allows it to recover full abstracts instead of empty placeholders. The agent also derives cleaner search
variants from a natural-language question by removing generic filler words, and it identifies itself to NCBI and
backs off on rate-limit responses, optionally using an API key to raise the request rate.

\section{Configuration and Scope}

Behaviour is centralized in the \texttt{Settings} dataclass in \path{config.py}, which holds the retrieval depth,
the coverage and support thresholds, the confidence floor and ceiling, the source-priority bonuses, the model
configuration, and the keyword lists used by the scope guard. Settings can be constructed directly or built from
environment variables, which is how the interfaces enable live retrieval and the language model when the
appropriate keys are present. The scope guard in \path{scope.py} uses the non-research keyword list to decline
clearly non-research queries while admitting research questions from any academic domain, returning an explicit
out-of-scope message that names the offending topic when possible.

\section{Ingestion of User Papers}

User-provided papers enter the system through \path{ingestion.py}, which parses Markdown, plain-text, JSON, and PDF
inputs into \texttt{PaperDocument} records. JSON is the most reliable format because it carries explicit metadata;
Markdown and text files are parsed for front-matter or a leading abstract paragraph; and PDFs are handled with a
best-effort text extractor. This is the mechanism by which the highest-priority evidence layer --- the user's own
papers --- is brought into the retrieval pool.

\section{Baselines as First-Class Clients}

The single-agent and RAG baselines in \path{baselines.py} are implemented as disciplined clients of the same
retrieval and scoring code rather than as separate systems. The single-agent baseline retrieves the top paper and
trusts a direct overlap check, while the RAG baseline retrieves several papers and aggregates simple supporting
overlap without explicit contradiction analysis; both cap their confidence to reflect their weaker evidential
discipline. Implementing the baselines on the shared substrate ensures that the comparison in the evaluation
isolates the effect of orchestration, claim-level verification, and critique, rather than confounding it with
unrelated differences in retrieval.

\section{Interfaces}

The Streamlit application in \path{app.py} exposes the full pipeline interactively: the user uploads papers, selects
a model, asks a question, and sees the executive summary, the retrieved papers with their source labels (uploaded,
corpus, or live API), the orchestration trace, and the per-claim verdicts with confidence and evidence. The FastAPI
service exposes the same capability programmatically. Both interfaces construct settings that enable live retrieval
and the language model from the environment, and both rely entirely on the package's reasoning logic, so the
interface and the studied system are one and the same.

\section{Implementation Principles}

Several principles recur across the codebase:

\begin{enumerate}
\item \textbf{Graceful degradation.} The language model and live retrieval are additive; the system retains full
deterministic functionality when they are unavailable.
\item \textbf{Grounding before generation.} Deterministic evidence scoring and verdict selection run first and
constrain any model refinement.
\item \textbf{Bounded effort.} Retrieval depth, claim counts, and the re-verification loop are bounded to keep the
workflow focused and predictable.
\item \textbf{Modularity.} Retrieval, extraction, verification, critique, synthesis, and reporting are separable
enough to be reused by the baselines and the interfaces.
\item \textbf{Traceability.} Each stage contributes explicit, typed outputs and trace lines that can be inspected
later by the user.
\end{enumerate}

These principles are what make AI-Scientist more than a prompt chain. They turn it into a controlled software
system whose behaviour can be studied systematically.

\section{Chapter Conclusion}

This chapter mapped the AI-Scientist architecture to its concrete implementation. The pipeline centers on a typed,
traceable set of agents behind a single Orchestrator; retrieval realizes the layered, source-prioritized evidence
model; verification and confidence are computed deterministically and refined optionally; the Critic drives bounded
adaptive re-verification; and the language model and live retrieval are optional enhancements that never override
the grounding rules. The same package drives the user interface, the API, and the baselines. With the
implementation now grounded in code, the next chapter describes how the system is evaluated.
